% Generated draft: requires mathematical and source review.
\documentclass[11pt]{article}
\usepackage[margin=25mm]{geometry}
\usepackage{fontspec}
\setmainfont{Latin Modern Roman}
\usepackage{amsmath,amssymb,mathtools,unicode-math}
\setmathfont{Latin Modern Math}
\usepackage{enumitem,needspace,xurl,hyperref}
\setlength{\parindent}{0pt}
\setlength{\parskip}{5pt}
\setlength{\emergencystretch}{3em}
\newcommand{\N}{\mathbb{N}}
\newcommand{\Z}{\mathbb{Z}}
\newcommand{\Q}{\mathbb{Q}}
\newcommand{\R}{\mathbb{R}}
\newcommand{\C}{\mathbb{C}}
\newcommand{\F}{\mathbb{F}}
\newcommand{\A}{\mathbb{A}}
\newcommand{\PP}{\mathbb{P}}
\newcommand{\E}{\mathbb{E}}
\newcommand{\eps}{\varepsilon}
\newcommand{\dd}{\,\mathrm{d}}
\begin{document}
\section*{UnsolvedMath ID: 700427}
\textbf{Source problem code:} catalogue entry
\subsection{Definitions and mathematical statement}
Let \(K\) be a commutative ring, let \(\Phi\) be a root system, and let \(G(K)\) be the Chevalley group associated with \(\Phi\) in the construction cited by the source, namely Steinberg's \emph{Lectures on Chevalley groups}. Write \(x_r(a)\) for a root element, and set
\[
x_r(B)=\{x_r(b):b\in B\}\qquad(B\subseteq K,\ r\in\Phi).
\]
The source specifies that \(G(K)=\langle x_r(K):r\in\Phi\rangle\).

For each \(r\in\Phi\), let \(\mathfrak A_r\) be an additive subgroup of \(K\). For every positive integer \(i\), define the set of powers
\[
\mathfrak A_r^i=\{a^i:a\in\mathfrak A_r\}.
\]
This notation denotes a set of powers, rather than an ideal power. The source calls \((\mathfrak A_r)_{r\in\Phi}\) an \emph{elementary carpet of type \(\Phi\) over \(K\)} when
\[
c_{ij,rs}\mathfrak A_r^i\mathfrak A_s^j
\subseteq\mathfrak A_{ir+js}
\qquad
\left(r,s,ir+js\in\Phi,\quad i,j\in\Z_{>0}\right),
\]
where \(c_{ij,rs}\) are the constants from the Chevalley commutator formula. Explicitly, each such inclusion requires
\[
c_{ij,rs}a^i b^j\in\mathfrak A_{ir+js}
\qquad\text{for every }a\in\mathfrak A_r\text{ and }b\in\mathfrak A_s,
\]
with the integer coefficient acting in \(K\).

\textbf{Unresolved qualification in the definition.} The displayed indexing is reproduced from the source. It does not exclude opposite roots \(s=-r\), whereas the usual Chevalley commutator formula defining these constants excludes opposite root pairs. For example, \(s=-r\), \(i=2\), and \(j=1\) give \(ir+js=r\in\Phi\), but do not supply a constant from that formula. The intended restriction or additional convention must therefore be clarified; no value for such a constant is imposed here. The precise Chevalley group construction and ring conventions also remain those of the cited reference, rather than additional hypotheses introduced in this draft.

\textbf{Problem, with these conventions fixed.} Determine necessary and sufficient conditions expressed in terms of the additive subgroups \((\mathfrak A_r)_{r\in\Phi}\) of an elementary carpet such that, on putting
\[
H(\mathfrak A)=\left\langle x_s(\mathfrak A_s):s\in\Phi\right\rangle\leq G(K),
\]
one has
\[
H(\mathfrak A)\cap x_r(K)=x_r(\mathfrak A_r)
\qquad\text{for every }r\in\Phi.
\]
The requested characterization is subject to resolution of the carpet-indexing qualification above.
\subsection{Short English statement}
Determine exactly which elementary carpets have the property that the subgroup generated by their prescribed root elements contains no additional elements in any root subgroup. The source's carpet condition leaves an unresolved convention for opposite roots. This is a statement-only draft requiring human mathematical review; the supplied registry labels the problem open.
\subsection{Sources}
\begin{itemize}
\item Supplied local UnsolvedMath registry record: permanent ID \(700427\), exact source problem code \texttt{catalogue entry}, title \emph{Kourovka Notebook Problem 7.28}. The supplied status is \texttt{open}; it is not an independently verified assessment.
\item V. M. Levchuk, proposer, \emph{The Kourovka Notebook (Unsolved Problems in Group Theory)}, original Problem 7.28, 7th issue (1980). The supplied record identifies the consulted edition as the 21st issue edition (2026), arXiv:1401.0300v48. Its collection description names E. I. Khukhro and V. D. Mazurov as editors. These bibliographic details are supplied attribution, not independently verified source consultation.
\item R. Steinberg, \emph{Lectures on Chevalley groups}, Yale Univ., New Haven, Conn., 1967; cited in the supplied statement for the construction of \(G(K)\).
\item The source additionally says ``See also 15.46.'' The text of that related problem was not supplied.
\end{itemize}
\end{document}
